% Portable, self-contained Pandoc publication template.
\documentclass[11pt,letterpaper]{article}
\usepackage[T1]{fontenc}
\usepackage{lmodern}
\usepackage{amsmath,amssymb,amsthm,mathtools}
\usepackage[margin=0.9in]{geometry}
\usepackage{microtype}
\usepackage{booktabs,longtable,array,calc,multirow}
\usepackage{etoolbox}
\usepackage{graphicx}
\usepackage{xurl}
\usepackage[numbers,sort&compress]{natbib}
\usepackage[hidelinks,unicode]{hyperref}
\usepackage{bookmark}
\urlstyle{same}
\setcounter{secnumdepth}{3}
\setlength{\emergencystretch}{3em}
\setlength{\parskip}{0.2em}
\allowdisplaybreaks[1]
\providecommand{\tightlist}{%
  \setlength{\itemsep}{0pt}\setlength{\parskip}{0pt}}
\providecommand{\passthrough}[1]{#1}
\providecommand{\pandocbounded}[1]{#1}
\AtBeginEnvironment{longtable}{\small\setlength{\tabcolsep}{4pt}}
\newenvironment{manualstatement}{%
  \par\addvspace{0.7\baselineskip}\begingroup\itshape\noindent\ignorespaces
}{\par\endgroup\addvspace{0.4\baselineskip}}
\newenvironment{manualdefinition}{%
  \par\addvspace{0.7\baselineskip}\begingroup\normalfont\noindent\ignorespaces
}{\par\endgroup\addvspace{0.4\baselineskip}}
\newenvironment{manualproof}{%
  \par\addvspace{0.4\baselineskip}\begingroup\normalfont\noindent\ignorespaces
}{\par\endgroup\addvspace{0.5\baselineskip}}
\title{What Remains Hard after ETH?\\[0.5em]{\large Exact Counting,
Better Representations, and a Refined Hardness Question}}
\author{Rohan Vinaik}
\date{\small 7 November 2026}
\begin{document}
\maketitle
\begin{abstract}
A subexponential algorithm can overturn ETH while leaving a very strong
form of computational hardness intact. The distinction is between
shaving logarithmic factors from an exponent and reducing that exponent
by a fixed power. This note explains the public successor hypothesis for
exact counting, why the task and its representation matter, and what a
Lean-verified finite tradeoff establishes for one algebraic counting
construction. It accompanies the technical article,
\emph{Refining Hardness after ETH}.
\end{abstract}
\section{A boundary can fail without the whole landscape
collapsing}\label{a-boundary-can-fail-without-the-whole-landscape-collapsing}

The Exponential Time Hypothesis made a precise prediction about
satisfiability: no deterministic algorithm would solve 3-SAT in
\(2^{o(n)}\) time, apart from polynomial input factors, where \(n\) is
the number of variables. The recent OpenAI and Claude papers report
algorithms that cross this boundary, including algorithms that compute
the exact number of satisfying assignments \citep{O26, Ant26}. If those
analyses are correct, the original ETH is false.

But the boundary was specific. Consider the difference between
\[2^{n/\sqrt{\log n}}\qquad\text{and}\qquad2^{\sqrt n}.\] Both are
subexponential. Their gains are very different. In the first, the
exponent remains larger than \(n^{0.99}\), or indeed
\(n^{1-\varepsilon}\) for any fixed positive \(\varepsilon\), once \(n\)
is sufficiently large. In the second, its power has dropped to one half.

That leaves a strong question after ETH: can exact counting achieve a
fixed-power improvement in its exponent, or does its worst-case
complexity retain the scale \(2^{n^{1-o(1)}}\)?

This question already has a public mathematical formulation. The
counting Satisfiability Hypothesis, \#SH, excludes exact CNF-counting
algorithms with running time \(O(2^{L^{1-\varepsilon}})\) for every
fixed \(\varepsilon>0\), where \(L\) counts the formula's variables,
clauses, and literal occurrences. MIT Hardness Group proves that it is
equivalent to Weak \#ETH, the corresponding variable-based hypothesis
for \#3-SAT with polynomial input factors \citep{SH26}. The hypothesis
remains unproved. The reported logarithmic improvements in the exponent
do not, by themselves, contradict it.

The technical article accompanying this note asks what structural
reasoning could support such a refined hardness claim \citep{Main26}.
Its aim is to identify the work an exact computation still has to do
after a better representation has delivered a substantial gain.

\section{What does the answer require us to
distinguish?}\label{what-does-the-answer-require-us-to-distinguish}

An exact count depends on which assignments satisfy a formula. It does
not follow that every assignment must be checked separately. Algebraic
summation, decomposition, and symmetry can combine many contributions in
one computation. Any argument for hardness has to account for those
joint operations.

A small example makes the role of access clear. Suppose we can inspect a
Boolean function only by asking whether a chosen assignment satisfies
it. If exactly one assignment satisfies the function, seeing that
positive answer certifies satisfiability immediately. It does not
certify an exact count of one. Every uninspected assignment might be
another satisfying assignment.

This difference persists even when all the alternative functions have
short width-two CNF descriptions. A singleton satisfying set needs only
unit clauses. A two-point satisfying set can be expressed with fixed
coordinates and signed equivalences. Under the declared point
observations, exact counting needs all \(2^n\) answers, while
satisfiability at a singleton needs one.

The example establishes an information requirement for those
observations. It does not establish exponential running time for a
program given the actual formula. Reading an explicit clause list
provides different information. In fact, reading all the bits of an
input acquires the whole input; the remaining difficulty lies in
computing with it.

This is the first lesson for a refined account of hardness. The task
determines which alternatives matter. The permitted operations determine
which alternatives can be resolved together. Their implementation
determines the computational cost. Counting candidates, or counting
certificates alone, does not settle all three questions.

A change of representation can therefore be a real algorithmic advance
even when it preserves every distinction relevant to the answer. It may
simply make those distinctions much cheaper to resolve jointly.
Transferring an old lower bound requires a stated relation between the
old and new operations, including their costs. The technical article
proves one finite observation-simulation criterion, and states its
limits explicitly.

\section{Where the algebraic gain
occurs}\label{where-the-algebraic-gain-occurs}

In one route through the OpenAI construction, a formula is converted
into a polynomial \(H\) whose value at zero is the exact count. The
algorithm computes Taylor data at other points and reconstructs the
required value. These data do not arrive for free. They are obtained
through measurements, substitutions, and sums over variable groups
\citep{O26}.

The reconstruction has a familiar information threshold. For arbitrary
rational polynomials of degree at most \(d\), initial Taylor
coefficients at distinct nonzero points determine the value at zero
exactly when their total number exceeds \(d\). Below that threshold,
adding the product of the required vanishing factors preserves all the
observed coefficients while changing the value at zero.

That statement still does not determine the algorithm's running time.
Nor does its necessity automatically hold within the narrower family of
polynomials generated by formulas: the added polynomial may not belong
to that family. A useful lower-bound argument must reach the actual
computational construction.

The Claude method provides another route. It rewrites the counting
problem as a network sum and uses exact identities to replace difficult
crossings by sums of simpler pieces. Derived table values are evaluated
by procedures, with their costs tracked recursively \citep{Ant26}. Its
gain comes from how contributions are combined. An argument assuming
that those procedures only inspect individual satisfying assignments
would miss the mechanism.

The common question is what interaction remains expensive after the
joint operation has been exploited. Recent analyses already address this
question within specified methods: Dias and Lima study interpolation
evaluation charges, Majdoddin studies restricted network recursions, and
Tulsyan studies support-constrained counting templates
\citep{DL26, Maj26, T26}. Their permitted operations and cost models
differ. The companion article compares them and supplies a fully checked
finite analysis and measurement application for the swap and
polynomial-fixing family. Its broader question remains how to make such
structural explanations precise across representations.

\section{A proved obstruction at the refined
scale}\label{a-proved-obstruction-at-the-refined-scale}

The construction has a finite space indexed by histories of coordinate
swaps. A polynomial in the swap-count operator is used to fix a large
subspace. Measurements control the remaining directions. A smaller
polynomial degree would reduce the substitutions needed later, but the
measurement budget limits how many unfixed directions can be handled.

The theorem proves this tradeoff using the actual operators. Their
spectrum has an explicit binomial distribution. A small number of
measurement rows forces the fixing polynomial to agree with alternating
signs on most of that distribution. Root counting and a binomial atom
bound then force the polynomial's degree to be at least a constant times
the square root of the number of swap positions.

The degree floor feeds into the specified recovery procedure. Write
\(n\) for its variable budget, \(R\) for the fixing polynomial's
positive degree budget, and \(Q\ge1\) for the number of swap positions.
The recovery uses \(q=\lfloor3nR/Q\rfloor\). Its complete substitution
family has size \[A(n,q)=\sum_{j=0}^{\min(q,n)}2^j\binom nj\] Under the
fixing condition and the construction's row and dimension budgets, the
theorem gives \[A(n,q)\ge
2^{\,n\log_2\log_2 n/(4\sqrt{\log_2 n})}
\qquad(n\ge2^{16}).\] It follows that for every fixed \(\varepsilon>0\),
beyond a threshold depending only on \(\varepsilon\),
\[2^{n^{1-\varepsilon}}\le A(n,q)\le3^n.\] The threshold is uniform over
every permitted fixing polynomial, measurement map, and parameter
choice. Thus changing the polynomial or measurement rows within these
contracts cannot give the complete enumeration a fixed-power improvement
in its exponent.

The source's sparse parity route uses different budgets and a different
Taylor cutoff. The same degree argument still forces a near-exponential
substitution family there, as the article shows \citep[Section
5.5]{Main26}. In both routes the result charges the routine that visits
each term. A faster way to compute the aggregate, or a different family
of operators, could change the cost.

The proof identifies a specific unresolved interaction and shows where
its cost reappears: measurement capacity forces agreement, agreement
forces degree, and degree forces enumeration. The spectral mechanism and
the near-exponential template scale have public antecedents. The Lean
development checks the finite chain and uniform eventual bound,
constructs the physical rows selected by the source algorithm, and
proves their fixing property, row count and height bounds
\citep{Main26}. The full formula-placement and counting algorithms
retain their separate source analyses.

\section{The broader claim remains a precise research
question}\label{the-broader-claim-remains-a-precise-research-question}

The unrestricted \#SH hypothesis concerns every deterministic
exact-counting algorithm. The theorem above concerns one defined family.
A proof of the former cannot be obtained just by removing the
qualifications from the latter.

What the result supplies is an example of the kind of account a broader
argument needs. It attaches hardness to task-changing alternatives and
the actual operations used to resolve them. It allows a representation
to produce a substantial gain, then proves a remaining tradeoff rather
than assuming that the original cost survives unchanged.

The reported ETH algorithms make this distinction necessary. They show
that an exponential-scale prediction can be too strong. They do not show
that exact counting admits a fixed-power reduction in its exponent.
Public successor hypotheses make that remaining claim legible and
testable; construction-specific theorems can explain particular reasons
it persists.

That is the direction developed in the companion paper: preserve the
serious hardness question, refine its quantitative statement, and prove
each structural obstruction at the scope its assumptions justify.
\bibliographystyle{plainnat}
\bibliography{references}
\end{document}
